\subsection{Measurement principles}
Table~\ref{tab:measures} separates person-level characteristics, session operations, proposed mediators and outcomes. Instruments are candidates; selection depends on population, language, permissions, recall period and psychometric performance. Brief bespoke state items require cognitive interviewing and pilot evaluation, and should not be described as validated scales. Trait alexithymia may be assessed with an authorised version of the TAS-20 \citep{bagby1994}; the Healthy-Unhealthy Music Scale \citep{hums2015} is an adjacent measure of music-use patterns, not a validated measure of AI-assisted creativity.

\begin{table}[htbp]
\centering\small
\caption{Candidate measurement architecture for AIM-ER.}
\label{tab:measures}
\begin{tabularx}{\linewidth}{@{}p{0.25\linewidth}p{0.19\linewidth}X@{}}
\toprule
Construct & Timing & Operational definition \\
\midrule
Goal and initial affect & Before session & Goal category, distress, valence and arousal; distinguish desired from felt emotion. \\
Trait characteristics & Baseline & Alexithymia, regulatory flexibility, musical expertise and prior AI experience. \\
Creative control & During session & Assigned interface; available and exercised editing, selection and revision operations. \\
Agency and ownership & After operations & Separate state ratings of perceived influence and felt belonging of the output. \\
Personal resonance & After exposure & Perceived correspondence between artefact and intended personal meaning. \\
Clarity and reappraisal & Before and after & Ability to identify emotions and change in interpretation; distinguish from liking. \\
Rumination and avoidance & After; follow-up & Repetitive negative thought and avoidance of relevant tasks or experiences. \\
Affective outcomes & Before, after, follow-up & Distress, valence and arousal on consistent scales; goal attainment separately. \\
Flexibility and functioning & Follow-up & Switching, stopping, return to intended activities; opportunity to switch recorded. \\
\bottomrule
\end{tabularx}
\end{table}
